%MB9T5
\setcounter{paragraph}{0}
\setcounter{ex}{0}
\PhanI
\loai{Sóng phản xạ}
\begin{ex}%[3P2Y8-1][Loai1]
Trong quá trình truyền sóng, khi sóng bị phản xạ. Tại điểm phản xạ thì sóng tới và sóng phản xạ sẽ
\choice
{luôn cùng pha}
{không cùng loại}
{luôn ngược pha}
{\True cùng tần số}
\loigiai{
Sóng tới và sóng phản xạ luôn cùng tần số.
}
\end{ex} \haidong{20}
\begin{ex}%[3P2B8-1][Loai1]
Một sóng nước lan truyền trên bề mặt nước tới một vách chắn cố định, thẳng đứng và phản xạ trở lại. Sóng tới và sóng phản xạ
\choice
{khác tần số, ngược pha}
{khác tần số, cùng pha}
{cùng tần số, cùng pha}
{\True cùng tần số, ngược pha}
\loigiai{
Sóng lan truyền gặp vật cản cố định bị phản xạ trở lại thì sóng tới và sóng phản xạ cùng tần số, ngược pha.
}
\end{ex} \haidong{20}
\begin{ex}%[3P2Y8-1][Loai1]
Khi nói về sự phản xạ của sóng cơ trên vật cản cố định, phát biểu nào sau đây đúng?
\choice
{Tần số của sóng phản xạ luôn lớn hơn tần số của sóng tới}
{\True Sóng phản xạ luôn ngược pha với sóng tới ở điểm phản xạ}
{Tần số của sóng phản xạ luôn nhỏ hơn tần số của sóng tới}
{Sóng phản xạ luôn cùng pha với sóng tới ở điểm phản xạ}
\loigiai{
\begin{itemize}
\item Tần số của sóng phản xạ luôn lớn hơn tần số của sóng tới $\Rightarrow$ Sai vì tần số của sóng phản xạ luôn bằng tần số của sóng tới.
\item Sóng phản xạ luôn ngược pha với sóng tới ở điểm phản xạ $\Rightarrow$ Đúng.
\item Tần số của sóng phản xạ luôn nhỏ hơn tần số của sóng tới $\Rightarrow$ Sai vì tần số của sóng phản xạ luôn bằng tần số của sóng tới.
\item Sóng phản xạ luôn cùng pha với sóng tới ở điểm phản xạ $\Rightarrow$ Sai vì vật cản cố định thì sóng phản xạ luôn ngược pha với sóng tới ở điểm phản xạ.
\end{itemize}
}
\end{ex} \haidong{20}
\loai{Điều kiện sóng dừng}
\begin{ex}%[3P2B8-1][Loai1]
Khi có sóng dừng trên một đoạn dây đàn hồi với hai điểm A, B trên dây là các nút sóng thì chiều dài AB
\choice
{bằng một phần tư bước sóng}
{bằng một bước sóng}
{bằng một số nguyên lẻ của phần tư bước sóng}
{\True bằng số nguyên lần nửa bước sóng}
\loigiai{
Điều kiện sóng dừng với hai đầu cố định (hai đầu là hai nút): $\ell =k\dfrac{\lambda }{2}\ ({k=1;2;3})$}
\end{ex} \haidong{20}
\begin{ex}%[3P2B8-1][Loai1]
Để có sóng dừng xảy ra trên một sợi dây đàn hồi với hai đầu dây có một đầu cố định và một đầu tự do thì chiều dài của dây phải bằng
\choice
{một số nguyên lần bước sóng}
{một số nguyên lần phần tư bước sóng}
{một số nguyên lần nửa bước sóng}
{\True một số lẻ lần một phần tư bước sóng}
\loigiai{
Để có sóng dừng xảy ra trên một sợi dây đàn hồi với hai đầu dây có một đầu cố định và một đầu tự do thì: 
$\ell = \left(k + \dfrac{1}{2}\right)\dfrac{\lambda}{2} = \dfrac{k\lambda}{2} + \dfrac{\lambda}{4}$.
}
\end{ex} \haidong{20}
\begin{ex}%[3P2Y8-1][Loai1]
Trên một sợi dây có chiều dài $\ell$, hai đầu cố định, đang có sóng dừng. Trên dây có một bụng sóng. Biết tốc độ truyền sóng trên dây là v không đổi. Tần số của sóng là
\choice
{$\dfrac{v}{\ell}$}
{\True $\dfrac{v}{2\ell}$}
{$\dfrac{2v}{\ell}$}
{$\dfrac{v}{4\ell}$}
\loigiai{
Sóng dừng hai đầu cố định với số bụng $n = 1$, do đó tần số: $f=n\dfrac{v}{2\ell}=\dfrac{v}{2\ell}$.
}
\end{ex} \haidong{20}
\loai{Khoảng cách bụng - nút}
\begin{ex}%[3P2Y8-1][Loai1] 
\nguon{QG - 2017} Một sợi dây căng ngang đang có sóng dừng. Sóng truyền trên dây có bước sóng $\lambda$. Khoảng cách giữa hai bụng liên tiếp là
\choice
{\True $\dfrac{\lambda}{2}$}
{$\dfrac{\lambda}{4}$}
{$2\lambda $}
{$\lambda $}
\loigiai{
Khoảng cách giữa hai bụng liên tiếp là $\dfrac{\lambda}{2}$.
}
\end{ex} \haidong{20}
\begin{ex}%[3P2B8-1][Loai1]
\nguon{MH - 2017} Trên một sợi dây đang có sóng dừng, sóng truyền trên dây có bước sóng là $\lambda$. Khoảng cách giữa hai nút sóng liên tiếp bằng
\choice
{$2\lambda $}
{\True $\dfrac{\lambda}{2}$}
{$\lambda $}
{$\dfrac{\lambda}{4}$}
\loigiai{
Trên sợi dây có sóng dừng khoảng cách giữa hai nút liên tiếp là một nửa bước sóng.
}
\end{ex} \haidong{20}
\begin{ex}%[3P2Y8-1][Loai1]
Trong hiện tượng sóng dừng trên dây. Khoảng cách giữa hai nút hay hai bụng sóng liên tiếp bằng
\choice
{một số nguyên lần bước sóng}
{một phần tư bước sóng}
{\True một nửa bước sóng}
{một bước sóng}
\loigiai{
Trong hiện tượng sóng dừng trên dây. Khoảng cách giữa hai nút hay hai bụng sóng liên tiếp bằng một nửa bước sóng.
}
\end{ex} \haidong{20}
\begin{ex}%[3P2Y8-1][Loai1]
Trên một sợi dây đàn hồi đang có sóng dừng. Khoảng cách từ một nút đến một bụng kề nó bằng
\choice
{một bước sóng}
{\True một phần tư bước sóng}
{hai bước sóng}
{một nửa bước sóng}
\loigiai{
Trên một sợi dây đàn hồi đang có sóng dừng. Khoảng cách từ một nút đến một bụng kề nó bằng một phần tư bước sóng.
}
\end{ex} \haidong{20}
\loai{Vị trí nút - bụng}
\begin{ex}%[3P2B8-1][Loai1]
Thực hiện sóng dừng trên dây AB có chiều dài $\ell$ với đầu B cố định, đầu A dao động theo phương trình $u=a\cos2\pi t$. Gọi M là điểm cách B một đoạn d, bước sóng là $\lambda$, và k là các số nguyên. Phát biểu nào sau đây là \textbf{sai}?
\choice
{Vị trí các nút sóng được xác định bởi công thức $d=\dfrac{k\lambda}{2}$}
{\True Vị trí các bụng sóng được xác định bởi công thức $d=(2k + 1)\dfrac{\lambda}{2}$}
{Khoảng cách giữa hai bụng sóng liên tiếp là $d=\dfrac{\lambda}{2}$}
{Khoảng cách giữa một nút sóng và một bụng sóng liên tiếp là $d=\dfrac{\lambda}{4}$}
\loigiai{
Vị trí các bụng sóng là $(2k + 1)\dfrac{\lambda}{4}$.
}
\end{ex} \haidong{20}
\loai{Nhiều phương án lựa chọn - Tìm phương án ĐÚNG}
\begin{ex}%[3P2B8-1][Loai1]
Phát biểu nào sau đây là đúng về hiện tượng sóng dừng? 
\choice
{Khoảng cách giữa hai bụng sóng là $\dfrac{\lambda}{2}$}
{Khi có sóng dừng trên dây có một đầu giới hạn tự do, điểm nguồn có thể là bụng sóng}
{\True Để có sóng dừng trên sợi dây đàn hồi với hai đầu là nút sóng thì chiều dài dây phải bằng một số nguyên lần nửa bước sóng}
{Khi có sóng dừng trên một sợi dây, hai điểm cách nhau $\dfrac{\lambda}{4}$ dao động vuông pha với nhau}
\loigiai{
\begin{itemize}
\item Đáp án C đúng vì khi xảy ra sóng dừng trên 2 đầu cố định thì $\ell = k\dfrac{\lambda}{2}$.
\item Đáp án A sai vì "thiếu từ gần nhất", vì khoảng cách giữa 2 bụng là $d = k\dfrac{\lambda}{2}$.
\item Đáp án B sai vì nguồn phải là nút sóng.
\item Đáp án D sai vì trong sóng dừng chỉ có cùng pha và ngược pha.
\end{itemize}
}
\end{ex} \haidong{20}
\begin{ex}%[3P2B8-1][Loai1]
Khi có sóng dừng trên dây đàn hồi thì
\choice
{nguồn phát sóng ngừng dao động còn các điểm trên dây vẫn dao động}
{\True trên dây có các điểm dao động mạnh xen kẽ với các điểm đứng yên}
{trên dây chỉ còn sóng phản xạ, còn sóng tới bị triệt tiêu}
{tất cả các điểm trên dây đều dừng lại không dao động}
\loigiai{
Khi có sóng dừng trên dây đàn hồi thì trên dây có các điểm dao động mạnh xen kẽ với các điểm đứng yên.
}
\end{ex} \haidong{20}
\loai{Thời gian giữa hai lần sợi dây duỗi thẳng}
\begin{ex}%[3P2K8-1][Loai1]
Một sợi dây chiều dài $\ell$ căng ngang, hai đầu cố định. Trên dây đang có sóng dừng với n bụng sóng, tốc độ truyền sóng trên dây là v. Khoảng thời gian giữa hai lần liên tiếp sợi dây duỗi thẳng là
\choice
{$\dfrac{v}{n\ell }$}
{$\dfrac{nv}{\ell }$}
{$\dfrac{\ell }{2nv}$}
{\True $\dfrac{\ell }{nv}$}
\loigiai{
$f=n\dfrac{v}{2\ell }\Rightarrow T=\dfrac{2\ell }{nv} \Rightarrow$ Khoảng thời gian giữa hai lần liên tiếp sợ dây duỗi thẳng là $\dfrac{T}{2}=\dfrac{\ell }{nv}$.
}
\end{ex} \haidong{20}
 
\begin{ex}%[3P2K8-1][Loai1]
Trong một thí nghiệm về sóng dừng trên một sợi dây có chiều dài $\ell$ với hai đầu cố định. Trên dây có sóng dừng với ba bụng sóng. Biết tốc độ truyền sóng trên dây là v. Thời gian giữa hai lần sợi dây duỗi thẳng liên tiếp là
\choice
{\True $\dfrac{\ell}{3v}$}
{$\dfrac{2\ell}{3v}$}
{$\dfrac{v}{3\ell}$}
{$\dfrac{2v}{3\ell}$}
\loigiai{
\begin{itemize}
\item $\lambda = \dfrac{2\ell}{3}\Rightarrow$ $T = \dfrac{\lambda}{v} = \dfrac{2\ell}{3v}$.
\item Thời gian giữa 2 lần sợi dây duỗi thẳng liên tiếp là:
$\dfrac{T}{2} = \dfrac{\ell}{3v}$.
\end{itemize}
}
\end{ex} \haidong{20}
\begin{ex}%[3P2K8-2][Loai1]
Cho A, B, C, D, E theo thứ tự là 5 nút liên tiếp trên một sợi dây có sóng dừng. M, N, P là các điểm bất kì của dây lần lượt nằm trong khoảng AB, BC, DE thì có thể rút ra kết luận nào sau đây?
\choice
{\True N dao động cùng pha P, ngược pha với M}
{M dao động cùng pha N, ngược pha với P}
{M dao động cùng pha P, ngược pha với N}
{Không thể kết luận được vì không biết chính xác vị trí các điểm M, N, P}
\loigiai{
\begin{itemize}
\item Trong sóng dừng, tất cả các điểm thuộc cùng 1 bó dao động cùng pha nhau và ngược pha với tất cả các điểm thuộc bó kế tiếp.
\item Vậy M và N ngược pha nhau, N và P cùng pha nhau.
\end{itemize}
}
\end{ex} \haidong{20}

\begin{ex}%[3P2B8-1][Loai1]
Khi có sóng dừng trên một sợi dây đàn hồi thì
\choice
{\True khoảng thời gian ngắn nhất giữa hai lần sợi dây duỗi thẳng là một nửa chu kì sóng}
{khoảng cách gần nhất giữa điểm nút và điểm bụng là một nửa bước sóng}
{khoảng cách gần nhất giữa điểm nút và điểm bụng là một bước sóng}
{tất cả các phần tử trên dây đều dừng lại (đứng yên)}
\loigiai{
Khi có sóng dừng trên một sợi dây đàn hồi thì khoảng thời gian ngắn nhất giữa hai lần sợi dây duỗi thẳng là một nửa chu kì sóng.
}
\end{ex} \haidong{20}

\loai{Trạng thái các điểm trên dây}

\begin{ex}%[3P2B8-1][Loai1]
Phát biểu nào sau đây là \textbf{sai} khi nói về sóng dừng xảy ra trên sợi dây?
\choice
{Khoảng cách giữa điểm nút và điểm bụng liền kề là một phần tư bước sóng}
{\True Hai điểm đối xứng với nhau qua điểm nút luôn dao động cùng pha}
{Khoảng thời gian giữa hai lần sợi dây duỗi thẳng là nửa chu kì}
{Khi xảy ra sóng dừng không có sự truyền năng lượng}
\loigiai{
Hai điểm đối xứng nhau qua nút luôn dao động ngược pha}
\end{ex} \haidong{20}

\PhanII
\setcounter{ex}{0}
\begin{ex}
Trên một sợi dây căng ngang, hai sóng cơ lan truyền ngược chiều giao thoa với nhau, tạo nên mô hình sóng dừng đặc trưng gồm các nút đứng yên và các bụng dao động.
\choiceTF[t]
{\True Sóng dừng là sóng được tạo thành do sự giao thoa của hai sóng tới và sóng phản xạ}
{\True Khi trên một sợi dây có sóng dừng, quan sát thấy trên dây có các bụng sóng xen kẽ với các nút sóng}
{Tại điểm phản xạ thì sóng phản xạ luôn ngược pha với sóng tới}
{Tại đầu tự do của dây thì sóng tới và sóng phản xạ ngược pha nhau} 
\loigiai{
\begin{itemchoice}
\itemch
\itemch
\itemch
\itemch
\end{itemchoice}
}
\end{ex} \haidong{20}
\begin{ex}
Trên một sợi dây đàn hồi căng ngang giữa hai điểm cố định, hai sóng cơ có cùng biên độ và tần số truyền ngược chiều nhau, tạo nên mô hình sóng dừng với những điểm nút và điểm bụng.
\choiceTF[t]
{\True Điểm bụng là điểm mà sóng tới và sóng phản xạ cùng pha}
{\True Điểm nút là điểm mà sóng tới và sóng phản xạ ngược pha}
{Trong sóng dừng có sự truyền pha từ điểm này sang điểm khác}
{\True Các điểm nằm trên cùng một bó thì dao động cùng pha}
\loigiai{
\begin{itemchoice}
\itemch
\itemch
\itemch
\itemch
\end{itemchoice}
}
\end{ex} \haidong{20}
\begin{ex}
Khi hai sóng cơ cùng tần số truyền ngược chiều trên cùng một môi trường gặp nhau, chúng giao thoa và tạo nên hiện tượng sóng dừng với các nút (điểm đứng yên) và bụng (điểm dao động cực đại).
\choiceTF[t]
{\True Sóng dừng có sự lan truyền dao động}
{\True Sóng dừng trên dây đàn là sóng ngang, trong cột khí của ống sáo, kèn là sóng dọc}
{\True Mọi điểm giữa hai nút liên tiếp của sóng dừng có cùng pha dao động}
{Bụng và nút sóng dịch chuyển với vận tốc bằng vận tốc lan truyền sóng}
\loigiai{
\begin{itemchoice}
\itemch
\itemch
\itemch
\itemch
\end{itemchoice}
}
\end{ex} \haidong{20}
\begin{ex}
Trên sợi dây căng theo phương thẳng đứng hai đầu cố định đang có sóng dừng.
\choiceTF[t]
{Không tồn tại thời điểm mà sợi dây duỗi thẳng}
{Trên dây có thể tồn tại hai điểm mà dao động tại hai điểm đó lệch pha nhau một góc là $\dfrac{\pi}{3}$}
{\True Hai điểm trên dây đối xứng nhau qua một nút sóng thì dao động ngược pha nhau}
{\True Khi giữ nguyên các điều kiện khác nhưng thả tự do đầu dưới thì không có sóng dừng ổn định} 
\loigiai{
\begin{itemchoice}
\itemch
\itemch
\itemch
\itemch
\end{itemchoice}
}
\end{ex} \haidong{20}

